\documentclass[12pt]{beamer}
\usepackage{beamerthemeCopenhagen}
\usecolortheme[rgb={0.18,0.3,0.3}]{structure}
%\usepackage[T1]{fontenc}
%\usepackage[utf8]{inputenc}
\usepackage[latin1]{inputenc}
\usepackage[brazil]{babel}
\usepackage{amsmath}
\usepackage{amsfonts}
\usepackage{amsthm}
\usepackage{pgf,tikz}
\usepackage{bbm}

\usetikzlibrary{arrows}

\newcommand{\caixa}{$\Box$}
\newtheorem{defi}{Defini��o}
\newtheorem{ex}{Exemplo}
\newtheorem{teo}{Teorema}
\newtheorem{cor}{Corol�rio}
\newtheorem{lema}{Lema}
\newtheorem{dist}{Distribui��o}
\newtheorem{emv}{M�xima verossimilhan�a}
\newtheorem{mo1}{Primeiro momento}
\newtheorem{mo2}{Segundo momento}
\newtheorem{prova}{Demonstra��o}
\newtheorem{BiNeg}{Defini��o}

\institute[IC-UFAL]{Universidade Federal de Alagoas\\
Instituto de Computa��o}
\title{Compara��o de estimadores pelo m�todo Monte Carlo das distribui��es: $\mathcal{G}_{\rm I}^ {0}$, Uniforme, Binomial negativa, Gamma e Exponencial}
\author{Ant�nio Marcos Larangeiras\and Alisson Nascimento\and Paulo\and Bruno\and Juliana}
\date{{\tiny Macei�-AL, Dezembro de 2013}}
\begin{document}
\begin{frame}
\maketitle
\begin{picture}(0,0)(0,-50)
\includegraphics[scale=0.11]{IC.png}
\end{picture}
\begin{picture}(0,0)(-250,-40)
\includegraphics[scale=0.13]{ufal.png}
\end{picture}
\end{frame} 

%=========================
\section{Introdu��o}
%=========================


%=======================================================================================================================================================
\begin{frame}
\frametitle{Introdu��o}
\begin{itemize}
\item Calcular o estimador por m�xima verossimilhan�a;

\bigskip

\item Calcular estimador pelo primeiro momento;

\bigskip

\item Calcular estimador pelo segundo momento;

\bigskip

\item Utilizar o M�todo Monte Carlo for�a bruta.

\end{itemize}
\end{frame}
%========================================================================================================================================================


%========================
\section{Distribui�\~oes}
%========================


%==============================================================================================
\subsection{Distribuição $\mathcal{G}_{\rm I}^{0}$}

\begin{frame}
\frametitle{Distribui\c c\~ ao $\mathcal{G}_{\rm I}^{0}$}

Denotaremos esta distribuição por $Z\sim\mathcal{G}_{\rm I}^{0}(z; \alpha,1,1)$, com $-\alpha, z>0$. Sua densidade é

\begin{displaymath}
f_{Z}=(z;\alpha,1,1)=\dfrac{\Gamma(L-\alpha)}{\Gamma(-\alpha)(1+z)^{1-\alpha}}, \mbox{com} -\alpha \ \mbox{e}\ z>0.
\end{displaymath}


\end{frame}

%========================================================================================================================================================
\begin{frame}
\frametitle{Distribui\c c\~ ao $\mathcal{G}_{\rm I}^{0}$}

A esperança matemática da distribuição $\mathcal{G}_{\rm I}^{0}$:

\begin{equation}
\label{eq:Gi0_esp}
E(Z)=\dfrac{\gamma}{\alpha-1}
\end{equation}\pause

A variância da distribuição $\mathcal{G}_{\rm I}^{0}$:

\begin{equation}
\label{eq:Gi0_var}
Var(Z)=\dfrac{1}{(\alpha-1)^2}+\dfrac{1}{(\alpha-2)(\alpha-1)^2}+\dfrac{1}{\alpha-1}
\end{equation}

\end{frame}

%====================
\subsubsection{Estimadores}
%====================
\begin{frame}
\frametitle{Estimadores da Distribuição $\mathcal{G}_{\rm I}^{0}$}
Pela Máxima Veressimilhança:
\begin{equation}
\label{eq:Gi0_MV}
\hat{\alpha}=\dfrac{n}{\displaystyle\sum_{i=1}^{n}\ln(1+z_i)}
\end{equation} \pause
Pelo momento central de ordem 2:
\begin{eqnarray}
\label{eq:Gi0_m_p2}
\hat{\alpha}_2=\overline{Z}
\end{eqnarray}
\end{frame}

\begin{frame}
\frametitle{Estimadores da Distribuição $\mathcal{G}_{\rm I}^{0}$}
Pelo segundo momento:
\begin{equation}
\label{eq:Gi0_m_p3}
\hat{\alpha}_3 = \frac{\frac{3}{k}\displaystyle\sum_{i=1}^{k}Z_{i}^{2}+2+\sqrt{\frac{1}{k^{2}}\displaystyle\sum_{i=1}^{k}Z_{i}^{4}+4}}{\frac{2}{k}\displaystyle\sum_{i=1}^{k}Z_{i}^{2}}
\end{equation}

\end{frame}


%===================
\subsubsection{Gráficos}
%=========================================================================================================================================================
\begin{frame}
\frametitle{Gráficos da distribuição $\mathcal{G}_{\rm I}^{0}$}

\end{frame}
%=========================================================================================================================================================




\subsection{Distribuição Uniforme}
%========================================================================================================================================================
\begin{frame}
\frametitle{Distribuição Uniforme $U(0, \theta)$}
\begin{dist}[$U(0, \theta)$]
Denotaremos esta distribuição por $X\sim U(x; 0,\theta)$, com $\theta>0$. Sua densidade é

{\small $$f(x) =
\left \{
\begin{array}{cc}
\frac{1}{ \theta } & 0 \leqslant x \leqslant \theta\\
0 & \textnormal{caso contrário.} \\
\end{array}
\right.$$ }

\bigskip\bigskip

{\scriptsize $E(X)=\dfrac{\theta}{2}$\hspace{1cm} e \hspace{1cm}$Var(X)=\dfrac{\theta^{2}}{12}$ }
\end{dist}
\end{frame}
%========================================================================================================================================================

%====================
\subsubsection{Estimadores}
%====================

\begin{frame}
\frametitle{Estimadores da distribuiç\~ao $U(0, \theta)$}
\begin{emv}[$U(0, \theta)$]
$$\widehat{\theta}_{1}={\dfrac{1}{\theta ^ n }}{\mathbb{1}_{\max{\{X_1,\dots,X_n\}\leqslant\theta}}},
\textnormal{portanto } \widehat{\theta} ={\max{\{X_1,\dots,X_n\}\leqslant\theta}} $$
\end{emv}

\begin{mo1}[$U(0, \theta)$]
$$\widehat{\theta}_{2}=2\overline{X}$$
\end{mo1} 

\end{frame}

\begin{frame}
\begin{mo2}[$U(0, \theta)$]
$$\widehat{\theta}_{3} = \sqrt{{{{3} \over {n}} {\sum_{i=1}^{n} X_{i}^{2}}}} $$
\end{mo2}
\end{frame}
%========================================================================================================================================================
%===================
\subsubsection{Gr\' aficos}
%===================
%=========================================================================================================================================================
\begin{frame}
\frametitle{Gráficos da distribuiç\~ao $U(0, \theta)$}

\end{frame}
%=========================================================================================================================================================

\subsection{Distribuição Binomial Negativa}
\begin{frame}
\frametitle{Distribuição Binomial Negativa}
Considere a situação de observar um evento dicotômico $Y_{i}$ independentes e identicamente distribuídos segundo uma lei de Bernoulli de probabilidade $p$. Suponha que se registre $Y$, o número de ensaios até obter exatamente $k$ sucessos. 

\end{frame}

\begin{frame}
\frametitle{Distribuição Binomial Negativa}

Seja uma variável aleatória que fornece o numero de ensaios até o k-ésimo sucesso. 
Assim, $Y$ tem uma distribuição binomial negativa com parâmetro $p \in (0,1)$, 
se sua função de probabilidade é dada por:

\begin{displaymath}
P_{r}(Y=y)=\left\{ \begin{array}{ll}
\binom{y-1}{k-1} \cdot p^{k} \cdot (1-p)^{y-k} & \textrm{se  y=k,k+1, \ldots.}\\
0 & \textrm{, caso contrário}
\end{array} \right.
\end{displaymath}


\end{frame}

%========================================================================================================================================================
\begin{frame}
\frametitle{Distribui\c c\~ ao $BN(p,k)$}

Usualmente sua função de probabilidade denota-se: $Y \sim BN(p,k)$.\pause

A esperança matemática da distribuição binomial negativa:

\begin{equation}
\label{eq:esp}
E(X) = {\frac{k}{p}}
\end{equation}\pause

A variância da distribuição binomail negativa:

\begin{equation}
\label{eq:var}
Var(X) = {\frac{k(1-p)}{p^{2}}}
\end{equation}

\end{frame}

%====================
\subsubsection{Estimadores}
%====================
\begin{frame}
\frametitle{Estimadores da Distribuição $BN(p,k)$}
Pela Máxima Veressimilhança:
\begin{equation}
\label{eq:BN_MV}
{\hat{p}} = {\frac{n k}{\sum_{i=1}^{n}{Y_{i}}}}
\end{equation} \pause
Pelo momento central de ordem 2:
\begin{eqnarray}
\label{eq:BN_m_p2}
{\hat{p}_{2}} = {\frac{-k + \sqrt{k^{2}+4k{Var(X_{i})}}}{2{Var(X_{i})}}}
\end{eqnarray}
\end{frame}

\begin{frame}
\frametitle{Estimadores da Distribuição $BN(p,k)$}
Pelo segundo momento:
\begin{equation}
\label{eq:BN_m_p3}
{\widehat{p}_{3}} = {\frac{{-k} + {\sqrt{k^{2} + 4 E(X_{i}^{2})(k^{2}+k)}}}{2E(X_{i}^{2})}}
\end{equation}

\end{frame}


%===================
\subsubsection{Gráficos}
%=========================================================================================================================================================
\begin{frame}
\frametitle{Gráficos da distribuição $BN(p,k)$}

\end{frame}
%=========================================================================================================================================================



%===================
\section{Resultados}
%===================

%===========================================================================================================================================================
\begin{frame}
\frametitle{Resultados}

\end{frame}
%============================================================================================================================================================

%=====================
\section{Conclus\~oes}
%=====================

%===========================================================================================================================================================
\begin{frame}
\frametitle{Conclus\~oes}

\end{frame}
%============================================================================================================================================================

\section{Refer�ncias Bibliogr�ficas}

\begin{frame}
\frametitle{Refer�ncias Bibliogr�ficas}


\bibliographystyle{plain}

\bibliography{Larangeiras}


\end{frame}
\end{document}